\subsection{Actual ordinary arithmetic seeds}
\label{sec:selmer-arithmetic-seed}

The arithmetic seed can be constructed at the original inertial type.
We use the prime-to-\(\ell\) character hypothesis of the actual
specialization argument; no uniform bound on the required weight is
asserted.

Fix the source's field and character data: \(\ell\) is odd,
\(K\subset L\) is imaginary quadratic inside the prescribed abelian
CM extension of \(\mathbb Q\), and \(\eta\) cuts out \(K\).
The characters \(\chi',\eta\) factor through
\(\mathbb Q(\zeta_N)\), with \(\ell\nmid N\), and take value
one at the prime \(w\nmid N\).  The nontrivial character \(\theta\)
has order \(\ell\), factors through \(\mathbb Q(\zeta_w)\),
and is trivial on \(G_L\).  For complex conjugation \(c\), assume
\[
 \ell\nmid\operatorname{ord}(\chi'),\qquad
 \chi'(c)=(-1)^j,\qquad
 \exists g\in G_L:\ \overline{(\epsilon^j\chi')(g)}^{\,2}\ne1,
 \quad \Psi=\epsilon^j\chi'.                         \tag{AS1}
\]
We use arithmetic Frobenius, so \(\epsilon(\operatorname{Frob}_q)=q\).

\begin{proposition}[The prescribed rank-two and rank-three seeds]
\label{prop:selmer-arithmetic-seed}
For every \(B\), the following constructions are possible on the
original field \(L\), with \(m\geq\max(B,2)\) and
\(m\equiv j\pmod{\ell-1}\).
If \(\chi'\ne\eta\), there is a continuous rank-three representation
\(r_3\), irreducible on every open subgroup of \(G_L\), with
residual characters \((\Psi,1,\Psi^{-1})\), exact inertia
\((\chi',1,\chi'^{-1})\) away from \(\ell\), and an ordered
ordinary flag of weights \((m,0,-m)\) and those residual characters
at every place above \(\ell\).

For the prescribed algebraic CM character \(\kappa\) of \(G_K\)
in the source's \(\mathrm{IsTwist}\) input, write \(\mathrm{ka}\)
for its prescribed restriction to \(G_L\).  There is a continuous
rank-two representation \(\sigma\) of \(G_{\mathbb Q}\) with
\[
 \det\sigma=\eta\theta,\qquad
 \overline\sigma^{\mathrm{ss}}=\overline{\Psi\eta}\oplus
       \overline\Psi^{-1}.
\]
Its restriction \(r_2\) to \(G_L\) is irreducible, has exact
away inertia \((\chi',\chi'^{-1})\), and ordered ordinary flag
of weights \((m,-m)\) with residual characters
\((\Psi,\Psi^{-1})\).

Both representations have a nondegenerate symmetric conjugate pairing
with multiplier one.  The representations \(r_3\epsilon^{-1}\)
and \(r_2\mathrm{ka}\), respectively, are attached in the original
normalization to regular algebraic conjugate self-dual cuspidal
automorphic representations.  Thus every field of the corresponding
\(\mathrm{IsSeed}\) record holds, including the residual
characteristic-polynomial identity at every Galois element.
\end{proposition}

\begin{proof}
\emph{An Eisenstein family at the prescribed conductor.}
Let \(\phi\ne1\) have prime-to-\(\ell\) order, conductor
\(M\mid N\), and \(\phi(w)=1\).  Then
\(\chi=\phi\theta\) is primitive of conductor \(C=Mw\),
where \(w\equiv1\pmod\ell\), \(w\geq7\), and
\(\ell\nmid C\).  Work over a finite coefficient DVR
\((\mathcal O,\lambda)\) containing the character and cusp values.
Choose \(a\) prime to \(C\ell\), with
\(a\equiv1\pmod{w\ell}\) and \(\overline{\phi(a)}\ne1\).
Then \(1-\phi(a)a^k\) is a unit for every \(k\).
The integral smoothed Kummer congruences
\cite[Main Theorem 8.2]{SeedDeligneRibet} give, for \(k\geq2\),
\[
 L(1-k,\chi)\equiv
 (1-\phi(w)w^{k-1})L(1-k,\phi)\equiv0\pmod\lambda.    \tag{AS2}
\]
Here the imprimitive \(\phi\) on conductor \(Mw\) removes the
\(w\)-Euler factor.  The same unit smoothing proves
\(L(1-k,\phi)\in\mathcal O\); the removed \(\ell\)-Euler
factors are one modulo \(\lambda\).  No division by \(k\) or a
Bernoulli denominator occurs, including when \(\ell=3\).

For \(\chi(-1)=(-1)^k\), put
\[
 E=\frac{L(1-k,\chi)}2+
       \sum_{n\geq1}\sum_{d\mid n}\chi(d)d^{k-1}q^n.
\]
The primitive constant-term formula
\cite[Definition 4.2 and Theorem 4.5(1)]{SeedDasguptaKakde}
vanishes away from the infinity class of \(\Gamma_0(C)\);
in that class it is a root-of-unity multiple of
\(L(1-k,\chi)/2\).  Hence every cusp constant on
\(\Gamma_1(C)\) lies in \(\lambda\mathcal O\).
The integral \(q\)-expansion principle
\cite[Corollary 1.6.2]{SeedKatz} makes \(E\) an integral form,
with first coefficient one.

\emph{A cusp lift with the same character, including small primes.}
Let \(X=X_1(C)_{\mathcal O}\), let \(D\) be its cusp divisor,
and let \(\pi:X\to Z=X/(\mathbb Z/C)^*\) be the coarse diamond
quotient.  For its Hodge line \(\omega\), form the coherent
\(\chi\)-eigensheaves
\[
 A_k=(\pi_*\omega^k)^\chi,\quad
 B_k=(\pi_*\omega^k(-D))^\chi,\quad
 C_k=(\pi_*(\omega^k|_D))^\chi.
\]
The sequence \(0\to B_k\to A_k\to C_k\to0\) is exact.
Indeed the target is supported at cusps, whose diamond stabilizers
have order dividing \(2C\): the cusp parabolic is generated by
\(-1\) and a unipotent with \(C\)-th power in \(\Gamma(C)\).
After strict henselization the local pushforward is induced from this
stabilizer.  Its order is invertible, so averaging over it proves
surjectivity.  This does not divide by the full diamond group order.

The invariant sections \(c_4^3,\Delta\) generate \(\omega^{12}\)
and induce the finite \(j\)-map.  Thus
\(\omega^{12}=\pi^*H\) for an ample line \(H\) on \(Z\),
with its given diamond linearization.  For each weight class modulo
12, relative Serre vanishing gives \(H^1(Z,B_k)=0\) for all
sufficiently large \(k\).  Lift the boundary of \(E/\lambda\)
to \(h\in H^0(A_k)\).  Then \(g=E-\lambda h\) is a cusp form
with unit first coefficient.  Its lattice is saturated in the modular
form lattice, since the boundary module is torsion free.  Therefore
\(g\bmod\lambda\) is a nonzero simultaneous Hecke eigenvector,
with eigenvalues \(1+\chi(q)q^{k-1}\) for \(q\nmid C\) and
\(U_q=1\) for \(q\mid C\).  This includes
\(T_\ell\equiv1\).
Deligne--Serre's eigenpacket-lifting lemma
\cite[Lemma 6.11]{SeedDeligneSerre} supplies a characteristic-zero
cuspidal eigenform after a finite coefficient extension.  Normalize it
to a newform \(f\).  Primitivity of \(\chi\) forces its exact
level to be \(C\mid Nw\); no auxiliary level was introduced.

\emph{The full local Galois conditions.}
Let \(\rho\) be the arithmetic-Frobenius Deligne representation
of \(f\).  Its determinant and residual semisimplification are
\[
 \det\rho=\epsilon^{k-1}\phi\theta,\qquad
 \overline\rho^{\mathrm{ss}}
       =1\oplus\overline{\epsilon^{k-1}\phi}.            \tag{AS3}
\]
Chebotarev extends the residual characteristic-polynomial identity
to every element.  At every \(q\mid C\), Carayol's local
compatibility theorem \cite[Theorem A]{SeedCarayol} applies over
\(\mathbb Q\).  Equality of the representation conductor and its
central-character conductor forces a principal series with one
unramified character.  A special or irreducible Weil representation
has strictly greater conductor than its determinant; for a principal
series the conductor is the sum of the character conductors, whereas
the product conductor is at most their maximum.  Hence monodromy
is zero and actual inertia is \(1\oplus\phi\theta\), including
at \(q=2\).

Kato's ordinary crystalline criterion, which makes no residual
irreducibility assumption \cite[Proposition 17.1 and Section 17.7]{SeedKato},
gives an unramified character \(\nu\) and the ordered flag
\[
              (\epsilon^{k-1}\phi\theta\nu^{-1},\nu).
                                                               \tag{AS4}
\]
Its unit-root eigenvalue is one modulo \(\lambda\), so
\(\nu\equiv1\) on the whole decomposition group.  The exact
inertia weights are \((k-1,0)\).
Irreducibility of \(\rho\) does not require irreducible residual
data: a split representation, by these local conditions and class
field theory, would consist of finite-order characters times
\(1,\epsilon^{k-1}\).  It would imply
\(|a_q(f)|\geq q^{k-1}-1\), contrary to the elementary cusp bound
\(a_n(f)=O_f(n^{k/2})\) for \(k>2\).

\emph{The rank-three representation.}
Take \(\phi=\chi'\eta\), \(k=m+1\), and
\[
 A=\operatorname{Sym}^2\rho\otimes(\det\rho)^{-1},
 \qquad r_3=A|_{G_L}.
\]
The parity condition holds, and \(m\) can be arbitrarily large
in its prescribed class.  Equations (AS3)--(AS4), with
\(\eta|_{G_L}=\theta|_{G_L}=1\), give the claimed characters
and ordered weights.  A nontrivial quadratic self-twist of \(\rho\)
would, on reducing (AS3), force
\(\overline{\epsilon^m\phi}\) to be quadratic.  This contradicts
(AS1).  The projective image is infinite by (AS4).
An irreducible two-dimensional representation with infinite projective
image and no quadratic self-twist has connected projective Zariski
closure \(\mathrm{PGL}_2\): the proper positive-dimensional cases
either fix a line or normalize a torus.  Every open subgroup has the
same connected closure, whose adjoint action is irreducible.

In the source's symmetric-square basis set
\[
 Q=\begin{pmatrix}0&0&-2\\0&1&0\\-2&0&0\end{pmatrix},
 \qquad P_3=QA(c).
\]
The displayed source matrix satisfies
\(S(g)^tQS(g)=(\det g)^2Q\).  Therefore \(P_3\) is symmetric
and invertible and
\(r_3(g)^tP_3r_3(cgc^{-1})=P_3\).

The Gelbart--Jacquet adjoint lift is cuspidal, since its quadratic
self-twist exception has been excluded
\cite[Theorem 9.3]{SeedGelbartJacquet}.
At a good prime with arithmetic Frobenius eigenvalues
\(\alpha_q,\beta_q\), the geometric Frobenius eigenvalues of
\(A\epsilon^{-1}\) are
\[
              q\{\alpha_q/\beta_q,1,\beta_q/\alpha_q\}.
\]
This is exactly the source's \(\operatorname{algSat}_3=q\operatorname{Sat}\).
Its infinitesimal character is \(\{m,0,-m\}\).
Iterate cyclic prime-degree base change along the original abelian
extension.  Arthur--Clozel's noncuspidal alternative would give a
nontrivial norm self-twist \cite[III.4.2 and III.5.1]{SeedArthurClozel}.
Dense Frobenius then gives a Galois self-twist by a character trivial
on the next subgroup.  Irreducibility on that subgroup and Schur's
lemma exclude it.  After the twist the conjugate pairing has multiplier
\(\epsilon^{-2}\); the two \(q\)-factors in
\(\operatorname{algSat}_3\) give conjugate self-duality by strong
multiplicity one.

\emph{The rank-two representation.}
Take \(\phi=\chi'^2\eta\), which is odd and nontrivial,
\(k=2m+1\), and
\[
            \sigma=\rho\epsilon^{-m}\chi'^{-1},
            \qquad r_2=\sigma|_{G_L}.
\]
Equations (AS3)--(AS4) give the exact determinant, residual characters,
inertia and ordered weights in the statement.
If \(\rho|_{G_L}\) were reducible, finite-index Clifford theory
would make it scalar, or give a quadratic self-twist trivial on
\(G_L\).  The scalar case contradicts infinite projective ordinary
inertia.  The self-twist case forces
\(\overline{\epsilon^{2m}\chi'^2\eta}|_{G_L}=1\), contrary
to (AS1).

Set
\(J=\left(\begin{smallmatrix}0&1\\-1&0\end{smallmatrix}\right)\)
and \(P_2=J\sigma(c)\).  We have \(\sigma(c)^2=1\),
\(\det\sigma(c)=-1\), and \(\det r_2=1\).
Consequently \(P_2^t=P_2\) and
\(r_2(g)^tP_2r_2(cgc^{-1})=P_2\).
Let \(\pi_f\) be the unitary representation of the newform, let
\(\Xi_\chi(v)=\chi'(\operatorname{Frob}_v)\) be the finite-order
Hecke character, and let \(\Xi\) be the prescribed Hecke avatar
of \(\mathrm{ka}\), using geometric Frobenius.  The required cusp
representation is
\[
 \Pi=\operatorname{BC}_{L/\mathbb Q}(\pi_f^\vee)
        \otimes|\det|^{1/2}\otimes\Xi_\chi\otimes\Xi.
\]
The same irreducibility argument preserves cuspidality at every
base-change step.  At a good place, the Satake factors from the first
two terms are \(q^{m-1/2}\{\alpha_q^{-1},\beta_q^{-1}\}\).
Multiplying by \(q^{1/2}\) in \(\operatorname{algSat}_2\)
therefore gives exactly the geometric Frobenius of
\(\rho\epsilon^{-m}\chi'^{-1}\mathrm{ka}\).
At a CM pair where \(\Xi\) has infinity types \((-1,0)\), the
infinitesimal characters are
\[
 \{m-\tfrac12,-m-\tfrac12\},\qquad
 \{m+\tfrac12,-m+\tfrac12\}.
\]
These are regular C-algebraic parameters and are negatives of each
other.  Finally \(\Xi\Xi^c=|\cdot|^{-1}\), equivalently the
Galois conjugate pairing multiplier is \(\epsilon^{-1}\).
The two \(q^{1/2}\) factors and strong multiplicity one give
conjugate self-duality.  The half-norm and the original
\(\mathrm{ka}\) have both been retained.
\end{proof}

This closes the exact seed input for the prime-to-\(\ell\) characters
used in the specialization theorem.  The prescribed CM-twist existence,
the marked point and prime-chain application, and the uniform integral
Selmer budget remain separate interfaces.  The construction introduces
no bound on coefficient degree or automorphic dimension.
